\chapter{Client API et cache}
\label{chap:api}

\section{Pagination de l'API JDM}

L'API JDM retourne au plus $n$ résultats par appel. Pour récupérer des listes complètes, \wido{} implémente une boucle de pagination native avec les paramètres \code{limit} et \code{offset}.

\begin{lstlisting}[caption=Boucle de pagination (extrait simplifié de jdmApi.js)]
let offset = 0, page = 0;
while (true) {
    if (page >= maxPages) break;         // Limite de sécurité
    if (results.length >= maxTotal) break;

    const url = `${base}?types_ids=${relId}
                 &limit=${limit}&offset=${offset}`;
    const data = await apiCall(url);

    results.push(...data.relations);
    page++; offset += limit;

    if (data.relations.length < limit) break; // Derniere page
}
\end{lstlisting}

Les paramètres de sécurité sont : \code{limit=1000}, \code{maxPages=5}, \code{maxTotal=5000}.

\section{Cache disque MD5}

Le cache disque évite les appels redondants à l'API JDM. Chaque clé de cache est un hash MD5 de l'URL ou des paramètres de la requête. Les données sont stockées en JSON sur le disque.

\begin{table}[H]
    \centering
    \begin{tabular}{ll}
        \toprule
        Type de cache & Clé \\
        \midrule
        Résolution de nœud & \code{node\_by\_name:\{name\}} \\
        Relations paginées & \code{relsPaged:\{dir\}:\{node\}:\{rel\}:...} \\
        Relations par ID & \code{relsPaged:\{dir\}:id:\{nodeId\}:\{rel\}:...} \\
        Cardinalité estimée & \code{cardinality:\{dir\}:\{node\}:\{rel\}:...} \\
        Vérification de couple & \code{check:\{id1\}:\{rel\}:\{id2\}} \\
        \bottomrule
    \end{tabular}
    \caption{Structure des clés de cache WIDO}
    \label{tab:cache}
\end{table}

\section{Gestion des erreurs API}

\begin{itemize}
    \item \textbf{HTTP 404} : nœud ou relation introuvable → résultat vide (cas légitime)
    \item \textbf{HTTP 500} : erreur serveur JDM → retry 1× après 500ms
    \item \textbf{Timeout ($>$15s)} : appel abandonné, warning remonté à l'interface
    \item \textbf{Erreur réseau} : log et continuation avec les données disponibles
\end{itemize}

La résilience du moteur repose sur ces mécanismes : une clause qui échoue ne bloque pas les autres.
